\setcounter{footnote}{59}

\printedpage{69}
% head: A. Grothendieck : Dérivateurs, Ch. VI, Édition des Universités de Montpellier II et Paris VII
which is given on objects by
\[\gamma(X) = C(X),\]
and which on arrows is defined in the obvious way in case A) (functoriality of $C(X)$ in $X$), and in the case of B) is constructed ad hoc, as the only functorial law for which the homomorphisms
\[i_X : X \longrightarrow C(X)\]
\msp{page 110}
(interpreted as arrows in $(\M/S)W_S^{-1}$) are themselves functorial in $X$.

We denote by
\[\Ho^{\mathrm{lc}}_W(S) \qquad\qquad \text{(lc as in ``locally constant'')}\]
the category of fractions $(\M/S)W_S^{-1}$, equivalent to $(\Fib S)W_S^{-1}$.

\circled{II} Let us take a localizer $W_S^0$ in $\M/S$
\[W_S^0 \subseteq W_S \subseteq \Fl(\M/S),\]
with the formal hypothesis similar to that on $W$ (cf. 1$^o$). We assume that

\begin{itemize}[leftmargin=4em]
\item[3$^o$)] $W_S^0 \cap \Fl(\Fib S) = W_S|\Fib S$.
\end{itemize}

(Cf. pages 94-96 for a general way of obtaining such localizers. One must clearly see that to define relative ``homotopy types'' over $S$ the localizer $W_S$ is much too coarse, morally it gives only a category of locally constant relative homotopy types, and this is precisely what is shown by the equivalence of \circled{I} , where $\Fib S$ is visualized as made up of the ``locally trivial fibrations'' over $S$.) We then have a second localized category % [sic] l'equivalence -> l'équivalence, accent manquant
\[\Ho_{W_S^0}(S) \overset{\text{def}}{=} (\M/S)(W_S^0)^{-1},\]
and two functors
\[
\begin{tikzcd}[column sep=large]
\Ho_{W_S^0}(S) \arrow[r, shift left, "\rho"] & \Ho^{\mathrm{lc}}_W(S), \arrow[l, shift left, "i"]
\end{tikzcd}
\tag*{$(*)$}
\]
where $\rho$ is the obvious (strict) localization functor, due to the inclusion $W_S^0 \subseteq W_S$, and where $i$ is
\msp{page 111}
essentially the functor deduced from the inclusion
\[(\Fib S, \underbrace{W_S|\Fib S}_{=W_S^0|\Fib S}) \longrightarrow (\M/S, W_S^0)\]
by localization \textit{by means of $W_S^0$}, hence the functor
\[(\Fib S)(W_S^0)^{-1} \longrightarrow (\M/S)(W_S^0)^{-1}, \tag*{$(**)$}\]
% foot: Version 9 novembre 2001
\printedpage{70}
% head: A. Grothendieck : Dérivateurs, Ch. VI, Édition des Universités de Montpellier II et Paris VII
taking into account the equivalence $(*)$ of \circled{I} , i.e. $i$ is the composite of the quasi-inverse of $(*)$, with the preceding functor $(**)$. Explicitly, it is defined by
\[i(X) = \text{class of C(X) in } (\M/S)(W_S^0)^{-1},\]
taking into account the fact that $X \longmapsto C(X)$ transforms $W_S$ into $W_S^o$, thanks to condition 3$^o$) on $W_S^0$. % [sic] W_S^o -> W_S^0, lettre o au lieu du chiffre 0 utilisé partout ailleurs

This being so, \textit{under hypothesis} B (``Quillen''), we have: in the pair of functors $(\rho, i)$ of $(*)$, page 110, \textit{$i$ is right adjoint to $\rho$,} i.e. we have for $\xi \in \mathrm{Ob}\,\Ho_{W_S^0}(S)$, $\eta \in \mathrm{Ob}\,\Ho^{\mathrm{lc}}_W(S)$
\[\Hom_{W_S^0}(\xi, i(\eta)) \simeq \Hom_{W_S}(\rho(\xi), \eta),\]
which is made explicit, for $X$, $Y$ in $\M/S$, by
\[\boxed{\Hom_{W_S^0}(X, C(Y)) \simeq \Hom_{W_S}(X, Y)}\]
(\footnote{\textbf{NB} The first member is independent of the choice of $W_S^0$! (Provided only that the hypothesis $W_S^0 \cap \Fl\,\Fib S = W_S \cap \Fl\,\Fib S$, in addition to $W_S^0 \subseteq W_S$, is satisfied.)}). It follows that \textit{the functor $i$ is fully faithful.}

The same adjunction formula, and the
\msp{page 112}
full faithfulness of $i$, hold in case A, when one strengthens hypothesis A (relative to a particular $S$) to a more restrictive and more comprehensive hypothesis, as follows.

\begin{itemize}[leftmargin=3em]
\item[A$'$)] The functor $C : \M/S \longrightarrow \M/S$ and the functorial morphism $i_X : X \longrightarrow C(X)$ are constructed ``in the manner of \textsc{Cartan-Serre}'' from a functor defined on the whole of $\M$
\[\M \xrightarrow{\mathrm{Ch}} \M,\]
of the ``paths'' type, i.e. equipped in addition with the following data
\begin{itemize}[leftmargin=3em]
\item[a)] \[X \xrightarrow{\alpha_X} \mathrm{Ch}(X) \xrightarrow{p_X = (\sigma_X, \beta_X)} X \times X\]
(structure of source and target of paths, and of trivial paths), and
\item[b)] A composition of paths, without requiring here the associativity of composition but only that trivial paths be two-sided units for composition. Moreover, we assume that for every $X$ % [sic] pout -> pour, faute de frappe
\[\begin{cases} \alpha_X \in W,\ p_X \in \mathrm{Fib} \\ (X \longrightarrow e) \in \mathrm{Fib}. \end{cases}\]
\end{itemize}
\end{itemize}
% foot: Version 9 novembre 2001

\printedpage{71}
% head: A. Grothendieck : Dérivateurs, Ch. VI, Édition des Universités de Montpellier II et Paris VII
\msp{page 113}
\textbf{Corollary.} Let $X$ be in $\M/S$, $Y$ in $\Fib S$, then
\[\Hom_{W_S^0}(X, Y) \xrightarrow{\sim} \Hom_{W_S}(X, Y).\]
\textsc{Proof}. By the adjunction formula (p. 111), the second member is $\Hom_{W_S^0}(X, CY)$. But as $Y$ and $C(Y)$ are in $\Fib S$, where $W_S$ and $W_S^0$ coincide, $Y \xrightarrow{i_Y} C(Y)$ is not only in $W_S$, but in $W_S^0$, hence $\Hom_{W_S^0}(X, Y) \xrightarrow{\sim} \Hom_{W_S^0}(X, CY)$, whence the result.

Let us return to Cat, Cat$/S$. The preceding digression on $(\M, W, \mathrm{Fib})$ has made it possible to better situate what concerns the functor $C$, and the relations between the localized categories of
\[\Cat/S \ \ (\text{for } W_S,\ W_S^{\mathrm{d}},\ W_S^{\mathrm{g}}) \qquad\quad \text{and} \qquad\quad \Parf S \ \ (\text{for } W_S^{\mathrm{f}} = W_S = W_S^{\mathrm{d}} = W_S^{\mathrm{g}})\]
as well as for the intermediate categories
\[\Parf S \subseteq \mathcal{F} \subseteq \Cat/S\]
full in Cat$/S$ (in the diagram p. 76, there are five others, namely $\Fibhot S$, $\Liss\,S$, $\Liss_0 S$, $\Prop\,S$, $\Prop_0\,S$). As for the relations between the localized categories of these seven categories, which are at the center of the diagram, the statements of theorem-scholie 1 (p. 70-81) give nothing more than the content of theorem-scholie 2 (p. 108-113). To establish them, we therefore have the choice between the two hypotheses A) or B) of the said theorem. Either use an explicit functor $C_S$, with the help of a
\msp{page 114}
functor $X \longmapsto \underline{\mathrm{Ch}}(X)$ (``paths in $X$'') on Cat, hence the structures of source, target and composition for paths, and trivial paths (acting as two-sided units for composition). The construction of such a functor $C$ does not look trivial and seems not to be known - I struggled to get there! (With, in addition, the associativity of composition \ldots). Or else use, in the manner of \textsc{Quillen}, the existence of a class of morphisms called Q-fibrant, with good factorization properties. This is not obvious in Cat either, and it is established by \textsc{Thomason}, by a somewhat complicated route. Perhaps I shall be led to come back to this - but only if necessary! % [sic] n'a par l'air -> n'a pas l'air, faute de frappe

On the other hand, the assertions concerning the categories $\Fib S$, $\Fibsc S$, or $\Cofib S$, $\Cofibsc S$, and their variants with indices 0 (thus the eight wing-categories of the diagram), and their relations to Cat$/S$ on the one hand, to $\Fibhot S$ on the other (for the categories with index 0), cannot be contained in a general carpet on categories of models $(\M, W)$, it seems to me, in view of the very special nature of the four categories designated above, of their very particular role in the categorical formalism. Now at least the ``split'' categories $\Fibsc S$, $\Cofibsc S$, and consequently also their variants with index 0, play an essential role for the theory of the HOT$(S)$, since it is they that
\msp{page 115}
intervene dually when we apply the general recipe for constructing a derivator from a category of models $(\M, W)$. So I see no simplification
% foot: Version 9 novembre 2001
\printedpage{72}
% head: A. Grothendieck : Dérivateurs, Ch. VI, Édition des Universités de Montpellier II et Paris VII
of substance to be made to theorem-scholie 1. At most one could prune the diagram of categories under consideration, by dropping $\Liss\,S$, $\Liss_0 S$, $\Prop\,S$, $\Prop_0 S$, which do not play a really useful role in it. As for the stated results, what concerns them is covered by what is said for the full subcategories $\mathcal{F}$ of Cat$/S$ which contain (say) $\Parf S$ (or at one's choice those which contain $\Fib S$ or $\Cofib S$).

The description of categories of the type HOT$(S)$, for fixed $S$, now seems to me perfectly well understood. It is time to let $S$ vary, and to understand the contravariant character in $S$. We thus take
\[S' \xrightarrow{f} S,\]
and we propose to describe the inverse image functors for the three types of categories HOT$(S)$. But we shall now restrict ourselves to the $\HOT_W(S)$ and $\HOT^{\mathrm{lc}}_W(S)$, since everything concerning $\HOT^o_W(S) = \HOT_W(S^o)$ \texttt{[rather} $\HOT_W(S^o)^o$\texttt{, cf. p. 50]} will follow from this by the duality $S \longmapsto S^o$. So it is a matter of describing
\msp{page 116}
\[f^* : \HOT_W(S) \longrightarrow \HOT_W(S') \tag*{$(*)$}\]
and the induced functor
\[f^* : \HOT^{\mathrm{lc}}_W(S) \longrightarrow \HOT^{\mathrm{lc}}_W(S').\]
When we describe the categories HOT$(S)$, in accordance with the general recipe, by means of functors
\[F : S^o \longrightarrow \Cat\]
(presheaves of categories on $S$), or, what amounts to the same thing, of split categories, the description of the arrow $(*)$ is obvious: it comes down to the law
\[\left|\begin{array}{rcl}
F & \longmapsto & F \circ f^o \\[1ex]
\underline{\mathrm{Hom}}(S^o, \Cat) & \longrightarrow & \underline{\mathrm{Hom}}(S'^o, \Cat),
\end{array}\right.\]
by passing to the localized categories for $W$-equivalence argument by argument. In the language of fibered categories, it is well known that this trivial operation on presheaves of categories corresponds to the \textit{inverse image} operation on categories over $S$ - which operation passes to splittings. This therefore gives us
\[f^*(\hot_{W,S}(X)) = \hot_{W,S'}(\underbrace{X \times_S S'}_{X'}),\]
for every $X$ split fibered over $S$. Now we know that the splitting of $X$ plays no role, in the end, in the description of the class
\[\hot_{W,S}(X) \in \HOT_W(S)\]
% foot: Version 9 novembre 2001
\printedpage{73}
% head: A. Grothendieck : Dérivateurs, Ch. VI, Édition des Universités de Montpellier II et Paris VII
\msp{page 117}
associated to $X$. We can forget the splitting, and even (to study the $\Hom_{\mathrm{Hot}_W(S)}(X, Y)$) work with $S$-functors \textit{of any kind} between categories fibered over $S$ - no need even to assume them cartesian, or even compatible with splittings. This seems to me a very important thing for the formalism of the $\HOT_W(S)$, and very special to the case of the category of models Cat. This is somewhat tacit in the theorem-scholie - but my initial doubts concerning the category $\underline{\mathrm{Fibsc}}$ itself show clearly that I felt there was something there that was not obvious, nor even (almost) \texttt{[?]} plausible. And yet the existence of the functor $\Phi$ shows that it is indeed so.

We can therefore interpret $f^*$ in terms of $\Fib S$, $\Fib S'$, by the base change functor
\[f^* : \Fib S \longrightarrow \Fib S', \qquad\quad X \longmapsto X \times_S S'/S',\]
which functor obviously preserves the fiberwise $W$-equivalence $W_S^{\mathrm{f}}$, $W_{S'}^{\mathrm{f}}$, which is the relevant one here
\[f^*(W_S^{\mathrm{f}}) \subseteq W_{S'}^{\mathrm{f}},\]
so that $f^*$ does indeed induce
\[(\Fib S)(W_S^{\mathrm{f}})^{-1} \longrightarrow (\Fib S')(W_S^{\mathrm{f}})^{-1}\]
\texttt{[rather} $(\Fib S')(W_{S'}^{\mathrm{f}})^{-1}$\texttt{]},
\msp{page 118}
i.e. the functor $(*)$
\[f^* : \HOT_W(S) \longrightarrow \HOT_W(S').\]
On the other hand, $\HOT_W(S)$ can also be described as a localized category of Cat$/S$ as a whole, with respect to the localizer $W_S^{\mathrm{d}}$, which on $\Fib S$, and even on $\Liss\,S$, coincides with $W_S^{\mathrm{f}}$. But here we run up against the difficulty that the inverse image functor
\[f^* : \Cat/S \longrightarrow \Cat/S', \qquad\quad X \longmapsto X \times_S S'/S',\]
is \textit{not} compatible in general (if $S$ is not isomorphic to the one-point category, for arbitrary $S'$) with the localizers $W_S^{\mathrm{d}}$, $W_{S'}^{\mathrm{d}}$. We see this e.g. by taking $S' = \{s\}$, and recalling that $W_S^{\mathrm{d}}$ does \textit{not} imply, on the whole of Cat$/S$, fiberwise $W$-equivalence.

These troubles do not occur on $\Liss\,S$, where $W_S^{\mathrm{d}}$ and $W_S^{\mathrm{f}}$ coincide. We can therefore describe $f^*$ on the HOT also with the help of % [sic] sûr -> sur, accent fautif
\[f^* : \Liss\,S \longrightarrow \Liss\,S', \qquad\quad X \longmapsto X \times_S S'/S',\]
by passing to the localized categories
\[f^* : (\Liss\,S)(W_S^{\mathrm{f}})^{-1} \longrightarrow (\Liss\,S')(W_{S'}^{\mathrm{f}})^{-1}.\]
In other words:

\textbf{Proposition 1.} \textit{If $X$ in} Cat$/S$ \textit{is smooth over $S$, then we have a canonical isomorphism}
\[f^*(\hot_{W,S}(X)) \simeq \hot_{W,S'}(X'), \qquad\quad \text{\textit{where} } X' = X \times_S S'/S'.\]
% foot: Version 9 novembre 2001